\label{sec:methods}
This section formalises rock--TBM interaction as an excavation-step-indexed
spatial entity-relation graph. The method starts from a TSP-derived rock voxel
field, an advancing parameterized TBM surface, and aligned monitoring records
under a common excavation-step index. It then constructs a sparse rock--TBM
interaction graph at each excavation step using active-zone, distance,
normal-compatibility, and excavation-state constraints. Finally, weighted
rock--machine edges are aggregated by TBM component to compute a component
$\times$ excavation-step anomalous interaction field. This field supports
excavation anomaly indication, component-indexed readout, and
jamming-risk-oriented interpretation while retaining traceable edge
contributions.
The graph is sparse because rock--machine edges are not fully connected; they
are retained only when the four spatial and excavation-state constraints are
jointly satisfied.
Figure~\ref{fig:pipeline} summarises the method structure.

The three subsections correspond to the input representation and two core
algorithmic operations. Section~\ref{sec:definition} defines the two graph
entity types and the aligned monitoring records. Section~\ref{sec:construction}
defines geometry-constrained sparse graph construction, and
Section~\ref{sec:interpretable_computation} defines component-level anomalous
interaction field computation from screened rock--machine edges.

The formulation estimates a queryable field rather than a direct
mechanical response. Given geological voxels $D_{\mathrm{geo}}$, a
parameterized TBM surface $M_{\mathrm{TBM}}$, and monitoring records aligned by
excavation step, the graph-construction operator
\[
\mathcal{G}: (D_{\mathrm{geo}},M_{\mathrm{TBM}},t)\mapsto G_t
\]
produces sparse rock--machine candidate relations, and the field-computation
operator
\[
\mathcal{F}: (G_t,\phi(g_i),\rho_j)\mapsto
\{I_c(t),\Delta I_c(t)\}_{c\in C,t=1,\ldots,T}
\]
aggregates those relations into a component-indexed anomalous interaction
field. The aligned monitoring records are used to read this field against
excavation responses; they are not inputs to $\mathcal{G}$ and do not define
graph edges.

\begin{figure}[htbp]
\centering
\includegraphics[width=0.95\linewidth]{fig1_method_framework.pdf}
\caption{Method overview of the geometry-constrained rock--TBM spatial
interaction graph computation. (a) Sparse graph construction from active rock
voxels and component-labelled TBM surface nodes using active-zone, distance,
normal-compatibility, and excavation-state constraints. (b) TSP-derived rock
voxel field with geological attribute support. (c) Parameterized TBM surface
components advanced along the excavation direction. (d) Component $\times$
excavation-step anomalous interaction field computed by aggregating weighted
rock--machine edges by TBM component; the field can subsequently be aligned
with monitoring responses and queried back to contributing edges.}
\label{fig:pipeline}
\end{figure}

\subsection{Rock, machine, and monitoring data representation}
\label{sec:definition}
The input data are organised into two graph entity types and aligned
monitoring records: rock voxels, TBM surface nodes, and excavation monitoring
responses aligned to excavation steps. Monitoring records are not graph nodes;
they are used after graph computation to evaluate and interpret the
component-level anomalous interaction field.

\paragraph{Rock voxel.} The geological input is a TSP-derived voxel field
$D_{\mathrm{geo}}=\{(c_i,g_i)\}$, where $c_i=(x_i,y_i,z_i)$ is the coordinate
of rock voxel $i$ and $g_i$ is its geological attribute vector (e.g.\
P-wave velocity $V_p$, S-wave velocity $V_s$, dynamic Young's modulus, velocity
ratio, Poisson's ratio, density). Rock voxels are the geological spatial
support of the graph.

\paragraph{TBM surface node.} The machine input is a parameterized TBM
surface. Each surface node is represented by $(p_j(0),n_j,\rho_j)$, where
$p_j(0)$ is the initial surface-node coordinate, $n_j$ is the local outward
surface normal, and $\rho_j$ is the machine component label. The current
implementation partitions the machine surface into cutterhead, front shield,
middle shield, and tail shield components; $\rho_j$ therefore indexes one of
these four labels. TBM surface nodes are the machine-side spatial support of
the graph.

\paragraph{Aligned monitoring records.} At each excavation step $t$ the monitoring
vector is
\[
u_t = [\mathrm{AdvanceRate}_t,\; \mathrm{RPM}_t,\; \mathrm{Torque}_t,\;
\mathrm{Thrust}_t,\; \mathrm{Penetration}_t,\; \mathrm{ShieldPressure}_t],
\]
and the response variables used downstream are denoted $r_t^{(k)}$, where $k$
indexes a response channel such as torque, thrust, penetration, advance rate,
or shield pressure. Monitoring responses are aligned to the graph output
$I_c(t)$ and $\Delta I_c(t)$; they do not participate in graph construction.

\paragraph{Common chainage coordinate.} Because TSP coordinates and TBM
monitoring mileage are commonly recorded in different local frames, the two
graph entity types and the monitoring records are mapped to a common local
chainage coordinate. With
$\mathrm{ShieldMileage}(t)$ the shield mileage stamp at step $t$, a simple
alignment is
\[
x_{\mathrm{local}}(t)=
\mathrm{ShieldMileage}(t)-\mathrm{ShieldMileage}_{\mathrm{start}},
\]
after which TBM surface nodes are advanced along the excavation direction
$d$:
\[
p_j(t+1)=p_j(t)+\Delta x\,d .
\]
Rock voxels, TBM surface nodes, and monitoring records are therefore aligned
to the same excavation-step index $t$. This subsection only establishes that
common index; it does not treat monitoring records as graph entities.

Table~\ref{tab:spatial_entities_relations} summarises the resulting
entity--relation data model. The table makes explicit that the graph is a queryable
entity-relation representation rather than only a plotting device or a
neural-network input tensor. Figure~\ref{fig:pipeline}(b--d) illustrates how
the rock voxel field, TBM surface components, and component-indexed field
readouts are formalised under the common excavation-step reference.

\begin{table}[htbp]
\centering
\caption{Entities, relations, and readouts in the geometry-constrained rock--TBM interaction
graph computation.}
\label{tab:spatial_entities_relations}
\scriptsize
\setlength{\tabcolsep}{3pt}
\begin{tabular}{p{0.17\linewidth}p{0.43\linewidth}p{0.30\linewidth}}
\toprule
Entity or quantity & Definition & Role \\
\midrule
$V_t^r$ & Active TSP-derived rock voxels & Geological spatial support \\
$V_t^m$ & Parameterized TBM surface samples & TBM component support \\
$E_t^{rr}$ & Rock--rock voxel neighbourhoods & Geological contiguity context \\
$E_t^{mm}$ & TBM surface neighbourhoods & Machine-surface topology context \\
$E_t^{rm}$ & Screened rock--machine candidate relations & Candidate interaction support \\
$A_c(t)$ & Sum of geometry weights for component $c$ & Auxiliary exposure readout \\
$I_c(t)$ & Geometry-weighted anomaly intensity for component $c$ & Main field readout \\
$e_{t+h}^{(k)}$ & Persistence residual of response $k$ & Aligned monitoring readout \\
\bottomrule
\end{tabular}
\end{table}

\subsection{Geometry-constrained sparse graph construction}
\label{sec:construction}
Algorithm~\ref{alg:graph_construction} summarises the graph construction
operation. Its purpose is to produce a sparse and queryable relation support
$E_t^{rm}$ for each excavation step, rather than to infer measured contact.

\begin{algorithm}[htbp]
\caption{Geometry-constrained sparse graph construction}
\label{alg:graph_construction}
\begin{algorithmic}[1]
\REQUIRE Rock voxels $\{(c_i,g_i)\}$, TBM surface nodes $\{p_j(0),n_j,\rho_j\}$, excavation steps $t=1,\ldots,T$
\ENSURE Graph snapshots $G_t=(V_t^r\cup V_t^m,E_t^{rr}\cup E_t^{mm}\cup E_t^{rm})$
\FOR{each excavation step $t$}
\STATE Advance the TBM surface to $p_j(t)$ using the common excavation-step coordinate.
\STATE Select active rock voxels $V_t^r$ using the face and shield active zones.
\STATE Construct $V_t^m$, $E_t^{rr}$, and $E_t^{mm}$ from voxel and TBM-surface neighbourhoods.
\FOR{each active rock voxel $i$ and TBM surface node $j$}
\IF{$d_{ij}(t)\leq \tau_{\mathrm{edge}}$, $\kappa_{ij}(t)\geq \eta_{\min}$, and $s_i(t)=1$}
\STATE Add $(i,j)$ to $E_t^{rm}$ and assign $w_{ij}^{rm}(t)=\exp[-d_{ij}(t)/\tau_{\mathrm{edge}}]\kappa_{ij}(t)$.
\ENDIF
\ENDFOR
\ENDFOR
\end{algorithmic}
\end{algorithm}

At each excavation step a graph snapshot is defined as
\[
G_t=\left(V_t^r \cup V_t^m,\; E_t^{rr} \cup E_t^{mm} \cup E_t^{rm}\right),
\]
where $V_t^r$ and $V_t^m$ are the rock and TBM surface node sets defined in
Section~\ref{sec:definition}, $E_t^{rr}$ denotes rock--rock voxel adjacency,
$E_t^{mm}$ denotes TBM surface adjacency, and $E_t^{rm}$ denotes candidate
rock--machine relations. Rock node features are
\[
x_i^r(t)=[c_i,g_i,s_i(t)],
\]
where $s_i(t)$ records whether voxel $i$ remains active and eligible for
interaction at step $t$. TBM node features are
\[
x_j^m(t)=[p_j(t),n_j(t),\rho_j].
\]
The intra-rock edge set $E_t^{rr}$ is constructed from voxel neighbourhoods,
and $E_t^{mm}$ is constructed from local neighbourhoods on the parameterized
TBM surface. The central graph-design decision is the construction of
$E_t^{rm}$, the rock--machine candidate edges.

\paragraph{Edge-screening rules.} A candidate edge $(i,j)$ is included in
$E_t^{rm}$ only when four conditions are jointly satisfied:
\[
\begin{cases}
c_i\in\Omega_t,\\
d_{ij}(t)\leq \tau_{\mathrm{edge}},\\
\kappa_{ij}(t)\geq \eta_{\min},\\
s_i(t)=1,
\end{cases}
\]
namely active-zone membership, a distance threshold, surface-normal
compatibility, and an excavation-state constraint. The first rule restricts
candidate relations to the excavation-relevant rock volume $\Omega_t$; the
second uses
\[
d_{ij}(t)=\|c_i-p_j(t)\|
\]
to keep only spatially close rock--machine pairs; the third uses the signed
surface-normal compatibility
\[
\kappa_{ij}(t)=
\max\left(0,\frac{n_j(t)^\top(c_i-p_j(t))}{d_{ij}(t)+\epsilon}\right)
\]
to keep only geometrically plausible facing relations; and the fourth prevents
already excavated voxels from repeatedly contributing to later graph snapshots.

The active zone $\Omega_t$ separates face-contact plausibility from
shield-side proximity within one active rock volume:
\[
\Omega_t=\Omega_t^{\mathrm{face}}\cup\Omega_t^{\mathrm{shield}} ,
\]
with, using cutterhead chainage $x_f(t)$, tail-shield chainage
$x_{\mathrm{tail}}(t)$, radial distance $r_i=\sqrt{y_i^2+z_i^2}$, cutterhead
radius $R_c$, shield radius $R_s$, face look-ahead length $L_f$, and radial
active-zone buffer $\tau_{\mathrm{zone}}$,
\[
\Omega_t^{\mathrm{face}}=
\left\{i \mid 0 \leq x_i-x_f(t) \leq L_f,\;
r_i \leq R_c+\tau_{\mathrm{zone}}\right\},
\]
\[
\Omega_t^{\mathrm{shield}}=
\left\{i \mid x_{\mathrm{tail}}(t)\leq x_i\leq x_f(t),\;
R_s \leq r_i \leq R_s+\tau_{\mathrm{zone}}\right\}.
\]
Table~\ref{tab:active_zone_parameters} lists the parameter values used in all
formal case runs. The radii and shield length come from the parameterized TBM
geometry. The edge-screening parameters $\tau_{\mathrm{edge}}$ and
$\eta_{\min}$ are varied in Section~\ref{sec:sensitivity_settings}. The
active-zone parameters $L_f$ and $\tau_{\mathrm{zone}}$ define the main
spatial support and are fixed in the reported formal runs.

\begin{table}[htbp]
\centering
\caption{Active-zone and candidate-relation parameters used in all formal
case runs.}
\label{tab:active_zone_parameters}
\scriptsize
\setlength{\tabcolsep}{5pt}
\begin{tabular}{lll}
\toprule
Parameter & Value & Role \\
\midrule
$L_f$ & 5.0\,m & Face look-ahead length for $\Omega_t^{\mathrm{face}}$ \\
$\tau_{\mathrm{zone}}$ & 5.0\,m & Radial active-zone buffer around face and shield \\
$R_c$ & 4.0\,m & Parameterized cutterhead radius \\
$R_s$ & 3.95\,m & Parameterized shield radius \\
Shield length & 10.0\,m & Axial support for $\Omega_t^{\mathrm{shield}}$ \\
$\tau_{\mathrm{edge}}$ & 2.0\,m & Candidate-edge distance threshold in the main setting \\
$\eta_{\min}$ & 0.3 & Minimum surface-normal compatibility in the main setting \\
Surface sampling & 0.5\,m & TBM surface-node sampling resolution \\
\bottomrule
\end{tabular}
\end{table}

\begin{figure}[htbp]
\centering
\includegraphics[width=0.95\linewidth]{fig3_geometry_constrained_edges.pdf}
\caption{Geometry-constrained rock--machine candidate relation construction.
Candidate relations are controlled by active-zone membership, distance decay,
surface-normal compatibility, and excavation state. The construction defines
spatially plausible relations; it does not recover measured contact.}
\label{fig:geometry_edges}
\end{figure}

\paragraph{Edge weight.} For each candidate relation the geometry-compatibility
weight is
\[
w_{ij}^{rm}(t)=
\exp\left(-\frac{d_{ij}(t)}{\tau_{\mathrm{edge}}}\right)\kappa_{ij}(t).
\]
The weight increases when a rock voxel is closer to the TBM surface and more
compatible with the local surface normal. It is not a contact force or
pressure. With $V_t^r$, $V_t^m$, $E_t^{rr}$, $E_t^{mm}$,
$E_t^{rm}$, and $w_{ij}^{rm}(t)$ defined, the graph entities of
Section~\ref{sec:definition} have become a computable rock--TBM interaction
graph.

\subsection{Component-indexed anomalous interaction field computation}
\label{sec:interpretable_computation}
Once the graph is constructed, the central graph computation is to aggregate
screened rock--machine edges into a component-level anomalous interaction
field. Response alignment and edge-level query are downstream readouts of this
field rather than separate method branches.

\begin{algorithm}[htbp]
\caption{Component-level anomalous interaction field computation}
\label{alg:interaction_field}
\begin{algorithmic}[1]
\REQUIRE Graph snapshots $G_t$, anomaly scores $q_i=\phi(g_i)$, component labels $\rho_j$
\ENSURE Component $\times$ excavation-step field $\{I_c(t),\Delta I_c(t)\}$
\FOR{each excavation step $t$}
\FOR{each TBM component $c$}
\STATE Select rock--machine candidate edges $E_{c,t}^{rm}=\{(i,j)\in E_t^{rm}\mid \rho_j=c\}$.
\STATE Compute geometric exposure $A_c(t)=\sum_{(i,j)\in E_{c,t}^{rm}}w_{ij}^{rm}(t)$.
\STATE Compute anomaly intensity $I_c(t)=\sum_{(i,j)\in E_{c,t}^{rm}}w_{ij}^{rm}(t)q_i/(A_c(t)+\epsilon)$.
\STATE Compute step change $\Delta I_c(t)=I_c(t)-I_c(t-1)$ when $t>1$.
\ENDFOR
\ENDFOR
\STATE Store edge-level contributions $s_{ij}(t)=w_{ij}^{rm}(t)q_i$ for subsequent query.
\end{algorithmic}
\end{algorithm}

\paragraph{Component-level field.} The geological contribution of each rock
voxel is represented by a restrained TSP-derived anomaly score
\[
q_i=\phi(g_i).
\]
In the main formulation $q_i$ is a bounded low-velocity anomaly score derived
from the training active-zone distribution of P-wave velocity:
\[
q_i=
\operatorname{clip}\left(
\frac{Q_{95}^{\mathrm{train}}(Vp)-Vp_i}
{Q_{95}^{\mathrm{train}}(Vp)-Q_{5}^{\mathrm{train}}(Vp)+\epsilon},
0,1
\right).
\]
This fixed training-reference scaling avoids per-chainage min--max
normalisation and keeps field values comparable along chainage; lower
$V_p$ values correspond to higher geological anomaly scores under a common
reference distribution.

For each TBM component $c$, let $V_c^m$ denote the set of TBM surface nodes
whose component label is $c$. Two component-level quantities are defined at
each excavation step.
The pure geometric exposure is
\[
A_c(t)=
\sum_{j\in V_c^m}
\sum_{i:(i,j)\in E_t^{rm}}
w_{ij}^{rm}(t),
\]
which reports the geometric exposure of component $c$ to the screened
candidate rock--machine relation set. The geology-weighted spatial interaction
intensity is
\[
I_c(t)=
\frac{
\sum_{j\in V_c^m}
\sum_{i:(i,j)\in E_t^{rm}}
w_{ij}^{rm}(t)q_i
}{
\sum_{j\in V_c^m}
\sum_{i:(i,j)\in E_t^{rm}}
w_{ij}^{rm}(t)+\epsilon
},
\]
which reports the weighted geological anomaly intensity in the screened
candidate relations incident to component $c$.
$I_c(t)$ is defined for cutterhead, front shield, middle shield, and tail
shield at each chainage step. Together with its step-level change
$\Delta I_c(t)$, it forms the component $\times$ excavation-step anomalous
interaction field. $A_c(t)$ is retained as an auxiliary exposure readout, not
as part of the anomalous interaction field itself. Here, the field denotes the
component $\times$ excavation-step matrix of $I_c(t)$ and $\Delta I_c(t)$
values, whereas a field readout denotes a downstream use of that matrix through
response alignment or edge-contribution query.

\paragraph{Step-level tracking.} Beyond the instantaneous intensity
$I_c(t)$, the field also supports a step-level change view
\[
\Delta I_c(t) = I_c(t) - I_c(t-1),
\]
which reports how the anomalous interaction intensity around component $c$
changes between consecutive excavation steps. Tracking $I_c(t)$ and
$\Delta I_c(t)$ over the test chainage interval gives a component-indexed
anomaly trajectory that indicates where and when the rock--machine
interaction intensity rises or falls along the alignment.

The computed field is subsequently aligned with monitoring residuals for
response readout and can be queried back to edge contributions. For response
alignment, short-horizon TBM responses are evaluated through persistence
residuals rather than raw responses:
\[
e_{t+h}^{(k)} = r_{t+h}^{(k)}-r_t^{(k)},
\]
where $h$ is the diagnostic horizon and $k$ indexes a response variable. The
diagnostic question is whether component-indexed field readouts co-vary with
the response change left after persistence has been removed:
\[
I_c(t)\rightarrow e_{t+h}^{(k)} .
\]
The reported association statistic is the Spearman correlation between
$I_c(t)$ and $e_{t+h}^{(k)}$ over the evaluated test chainage interval. This
readout links the graph-computed field to TBM response change; it is
interpreted as field--residual co-variation within the evaluated chainage
interval, not as causal attribution. The choice of diagnostic pairs, null
variants, detrending, and robustness checks is part of the case evaluation
design and is specified in Section~\ref{sec:diagnostic_readout_variants}.
For edge-level traceability, candidate edges for any fixed diagnostic sample
are ranked by their contribution
\[
w_{ij}^{rm}(t)\,q_i,
\]
and each contributing edge retains the rock voxel identity, the TBM surface
node identity, the component label $\rho_j$, the distance $d_{ij}(t)$, the
normal compatibility $\kappa_{ij}(t)$, the geometry weight $w_{ij}^{rm}(t)$,
the anomaly score $q_i$, and the weighted contribution. The graph therefore
not only produces a component-indexed field but also traces that field back to
specific rock--machine spatial relations.

\FloatBarrier
